\documentclass[11pt,letterpaper]{article}

% --- MÁRGENES Y ESPACIO ---
\usepackage[left=1.7cm,right=1.7cm,top=2cm,bottom=2.5cm, columnsep=0.75cm]{geometry}
\usepackage[utf8]{inputenc}

% --- MATEMÁTICA Y CIENCIA ---
\usepackage{breqn}
\usepackage[spanish, es-tabla]{babel}
\usepackage{amsmath}
\usepackage{amsfonts}
\usepackage{amssymb}
\usepackage{siunitx}
\sisetup{mode=text, output-decimal-marker = {.}, per-mode = symbol, qualifier-mode = phrase, qualifier-phrase = { de }, list-units = brackets, range-units = brackets, range-phrase = --}
\decimalpoint
\usepackage{cancel}

% --- TABLAS Y FIGURAS ---
\usepackage{graphicx}
\usepackage[margin=10pt,labelfont=bf]{caption}
\captionsetup[figure]{position=bottom}
\captionsetup[table]{position=bottom}
\usepackage{subcaption}
\usepackage{floatrow}
\usepackage{longtable}
\usepackage{array}
\usepackage{xtab}
\usepackage{multirow}
\usepackage{colortab}
\usepackage{colortbl}
\usepackage{tabularx}
\usepackage{stfloats}
\usepackage{booktabs}
\usepackage{float}
\floatstyle{plaintop}
\newfloat{anexo}{thp}{anx}
\floatname{anexo}{Anexo}
\restylefloat{anexo}

% --- FORMATO Y TIPOGRAFÍA ---
\usepackage{makeidx}
\usepackage{xcolor}
\usepackage[stable]{footmisc}
\usepackage[section]{placeins}
\usepackage{ulem}
\usepackage{lipsum}
\usepackage{multicol}
\usepackage{palatino}
\usepackage{titlesec}
\usepackage{enumitem}
\usepackage{comment}
\usepackage{dirtytalk}
\usepackage{microtype}

\titleformat*{\section}{\bfseries\large}
\titleformat*{\subsection}{\bfseries\normalsize}

% --- ENCABEZADOS Y PIES DE PÁGINA ---
\usepackage{fancyhdr}
\fancyhead[L]{Laboratorio de Sistemas Dinámicos} 
\fancyhead[R]{Braunstein - Prado Rodriguez - Informe Interno} % Actualizado para el informe actual
\renewcommand{\headrulewidth}{0.09pt} 
\fancyfoot[C]{\thepage} 
\pagestyle{fancy} 

\fancypagestyle{firstpage}{%
  \fancyhf{} 
  \renewcommand{\headrulewidth}{0pt} 
}

% --- ENLACES Y BIBLIOGRAFÍA ---
\usepackage{xurl} 
\usepackage[colorlinks=true, linkcolor=orange, urlcolor=blue, citecolor=red, anchorcolor=orange, breaklinks]{hyperref}
\usepackage{csquotes}
\usepackage[style=numeric, sorting=none, backend=biber]{biblatex}

\DeclareFieldFormat{url}{\href{#1}{[Enlace]}}
\DeclareFieldFormat{title}{#1} 
\DeclareFieldFormat[article]{title}{#1} 
\renewcommand{\bibfont}{\normalfont\normalsize} 

\addbibresource{referencias.bib}

%%%%%%%%%%%%%%%%%%%%%%
% Inicio del documento
%%%%%%%%%%%%%%%%%%%%%%
\raggedbottom 
\setlength{\columnsep}{20pt} 

\begin{document}

\thispagestyle{firstpage} 

\renewcommand{\labelitemi}{$\checkmark$}
\renewcommand{\CancelColor}{\color{red}}
\newcolumntype{L}[1]{>{\raggedright\let\newline\\\arraybackslash}m{#1}}
\newcolumntype{C}[1]{>{\centering\let\newline\\\arraybackslash}m{#1}}
\newcolumntype{R}[1]{>{\raggedleft\let\newline\\\arraybackslash}m{#1}}

\begin{center}
    % TÍTULO
    {\LARGE \textbf{Gran Manual del ÑANDÚ LSD}}\\
    \vspace{2mm}
    {\LARGE \textbf{Informe Interno de Laboratorio 6 y 7}}\\
    \vspace{4mm}
    
    % AUTORES
    {\large \textbf{Braunstein Lucas}, \textbf{Prado Rodriguez Santiago}}\\
    \vspace{2mm}
    {\small braunstein.lucas@gmail.com, santiagopradorodriguez@gmail.com}\\
    
    \vspace{4mm}
    {\large \textbf{Directora:} Ana Amador}\\
    \vspace{1mm}
    {\large \textbf{Colaborador:} Felipe Cignoli}\\
    
    \vspace{6mm}
    \includegraphics[width=0.5\linewidth]{insumos_ia/cuaderno_de_labo/images/logo.png}
    \vspace{8mm}
    
    % INFO DEL LABO
    {\large \textit{Laboratorio de Sistemas Dinámicos\\
    Departamento de Física, FCEyN, UBA\\
    \today}}
\end{center}

\newpage

\section*{Introducción}
\addcontentsline{toc}{section}{Introducción}

Este documento constituye el informe interno detallado sobre el diseño, construcción y caracterización del hardware utilizado en el proyecto. 
Se documenta específicamente la placa electrónica de amplificación de tres canales, el circuito de protección de polaridad, y la estructura de interconexión con el 
software de adquisición Ñandú LSD. El objetivo es servir como manual de usuario y registro técnico de las decisiones de laboratorio tomadas a lo largo del año. 
Además, compila y analiza todo el proceso experimental y computacional llevado a cabo durante los laboratorios para lograr la extracción y clasificación fonética 
de vocales a partir de señales de electromiografía de superficie (sEMG) de los músculos del tracto vocal.

Hay información del desarrollo de software y los algoritmos utilizados para la decodificación de las vocales del español. Hay procedimientos, mediciones y resultados.
 Se aconseja leer los informes elaborados (que son 3). Toda la información y mediciones se encuentran en un disco duro y la estructura sigue la idea del índice.

Lxs saludan Lucas y Santiago.

\newpage
{\hypersetup{linkcolor=black}
\tableofcontents}
\newpage

% --- INCLUSIÓN DE SECCIONES ---
% Antigravity/vos van a generar estos archivos por separado.

\subsection{Introducción}
El presente diseño se basa en el uso de amplificadores de instrumentación diferenciales, específicamente el circuito integrado AD620, ideal para el registro de señales bioeléctricas de baja amplitud. Como punto de partida, se utilizó y adaptó un circuito preamplificador previamente desarrollado por un grupo de Laboratorio 7, el cual había sido diseñado originalmente para registrar potenciales de activación en músculos de aves. 

Una de las adaptaciones más críticas consistió en el ajuste del filtro pasa-altos de entrada. El circuito original empleaba una frecuencia de corte de \SI{72}{\hertz}, adecuada para conservar la frecuencia fundamental de los músculos aviares. Sin embargo, dado que este proyecto se enfoca en el reconocimiento electromiográfico de músculos faciales humanos, cuyo contenido espectral de interés puede encontrarse en frecuencias más bajas, fue necesario readaptar el filtrado. Tras revisar bibliografía específica (como el trabajo de Ratnovsky), se determinó que filtrar frecuencias por debajo de \SI{50}{\hertz} resultaba más conveniente para conservar la información de los músculos faciales sin comprometer la señal.

\subsection{Diseño de circuito}
El circuito diseñado consta de tres canales de amplificación independientes y un sistema central de protección. Para garantizar la seguridad de los componentes, en especial los costosos amplificadores AD620, se implementó un circuito de protección contra inversión de polaridad. Este circuito utiliza transistores MOSFET tipo P (F9540N) y tipo N (IRFZ44N), complementados con fusibles de \SI{0.25}{\ampere}, evitando posibles daños en caso de conectar erróneamente las baterías de alimentación.

\begin{figure}[htbp]
\centering
\includegraphics[width=0.8\textwidth]{rutas/pendientes.png}
\caption{Circuito esquemático mostrando los canales de amplificación diferencial y la etapa de protección contra cambio de polaridad.}
\label{fig:circuito_esquematico}
\end{figure}

Un aspecto fundamental del diseño fue la determinación de la resistencia de ganancia ($R_g$) para los amplificadores AD620. La amplificación del integrado está dada por la relación de fábrica dictada por la hoja de datos. Durante la etapa experimental, se realizaron mediciones empíricas sobre el músculo facial \textit{depressor anguli oris} y barridos de frecuencia para caracterizar la respuesta del sistema. 

Inicialmente se probaron resistencias de bajo valor, como \SI{47}{\ohm}, buscando maximizar la ganancia (la cual apuntaba a valores superiores a \num{1000}). Sin embargo, los resultados mostraron que forzar el integrado a factores de amplificación muy elevados generaba saturación, distorsión y un aumento considerable del nivel de ruido en las señales registradas. Adicionalmente, los ensayos revelaron que operar en los límites de amplificación de este circuito integrado comprometía la linealidad y el ancho de banda del canal. 

Por estos motivos, se descartaron los valores bajos de resistencia y se optó definitivamente por utilizar resistencias $R_g$ de \SI{100}{\ohm} en los tres canales. Este valor proporciona una ganancia estable y segura (aproximadamente $495\times$), evitando la saturación del integrado frente a las tensiones entregadas por las baterías y manteniendo la amplificación dentro de un régimen lineal altamente confiable para la adquisición de la electromiografía.

Respecto al suministro de componentes críticos, los amplificadores de instrumentación AD620 fueron adquiridos en G.M. ELECTRÓNICA S.A., Av. Rivadavia 2458, C1034ACQ Cdad. Autónoma de Buenos Aires, 011 4953-0417.


\input{sections/02_software.tex}

\newpage
\vspace*{2cm}
\begin{center}
    \Huge \textbf{Capítulo 3: Experimentos}
\end{center}
\vspace{2cm}
\phantomsection
\addcontentsline{toc}{section}{Experimentos}

\subsection{Experimento de Labo 6}

Para la realización del experimento principal de Laboratorio 6, el objetivo fue registrar la actividad de tres músculos faciales involucrados en el habla, de acuerdo a la bibliografía consultada. Vamos a visualizar y ubicar los músculos que nos interesan según el paper de Ratnovsky (2021); apuntamos a poner electrodos en los siguientes músculos:

\begin{figure}[H]
    \centering
    \includegraphics[width=0.7\linewidth]{insumos_ia/cuaderno_de_labo/images/image252.png}
    \caption{Ubicación anatómica de los músculos faciales según Ratnovsky et al. (2021).}
    \label{fig:ratnovsky2021}
\end{figure}

\begin{figure}[H]
    \centering
    \includegraphics[width=0.6\linewidth]{insumos_ia/cuaderno_de_labo/images/image541.png}
    \caption{Detalle del músculo Orbicularis oris.}
    \label{fig:orbicularis}
\end{figure}

\begin{figure}[H]
    \centering
    \includegraphics[width=0.7\linewidth]{insumos_ia/cuaderno_de_labo/images/image513.png}
    \caption{Detalle anatómico del músculo Mylohyoid.}
    \label{fig:mylohyoid}
\end{figure}

Una vez definidas las posiciones anatómicas, procedimos a colocar los electrodos sobre el sujeto de prueba (Santiago). Para evitar confusiones respecto a versiones anteriores del montaje, redefinimos varias veces el nombre de los canales. Finalmente, preferimos nombrarlos según los \textit{inputs} del NIDAQ (la placa de adquisición). La asignación definitiva fue la siguiente:
\begin{itemize}
    \item \textbf{CH2:} Conectado al músculo \textit{Depressor anguli oris}.
    \item \textbf{CH1:} Conectado al músculo \textit{Orbicularis oris}.
    \item \textbf{CH0:} Conectado al músculo \textit{Mylohyoid}.
\end{itemize}

Comenzamos probando el \textbf{CH2} en el depresor. Aplicamos un filtro notch (\SI{50}{\hertz}) y la señal mejoró drásticamente; sin el filtro se observaba excesivo ruido de línea. 

A continuación, sumamos el \textbf{CH1} al orbicularis. Observamos si el sistema captaba ambas señales en simultáneo. Efectivamente, aunque se percibió un poco de ruido, la adquisición simultánea fue exitosa.

Finalmente, sumamos el \textbf{CH0} al mylohyoid. El sistema capturó los tres canales correctamente. Sin embargo, notamos un problema mecánico: el formato de los electrodos (\enquote{curitas}) que usamos debajo del mentón no era el óptimo. Al tener que luchar contra la gravedad y soportar el peso adicional de las fichas cocodrilo, requerían de un mejor agarre. Como solución práctica, procedimos a armar dos electrodos de los viejos, que ofrecían mayor adherencia y estabilidad mecánica.

A continuación, se muestran las imágenes del montaje final de los electrodos sobre el sujeto de prueba:

\begin{figure}[H]
    \centering
    \includegraphics[width=0.32\linewidth]{insumos_ia/cuaderno_de_labo/images/image97.jpg}
    \includegraphics[width=0.32\linewidth]{insumos_ia/cuaderno_de_labo/images/image241.jpg}
    \caption{Disposición frontal y lateral de los electrodos en el sujeto. (Imágenes extraídas del cuaderno de laboratorio).}
\end{figure}

\begin{figure}[H]
    \centering
    \includegraphics[width=0.45\linewidth]{insumos_ia/cuaderno_de_labo/images/image452.jpg}
    \includegraphics[width=0.45\linewidth]{insumos_ia/cuaderno_de_labo/images/image362.jpg}
    \caption{Detalle de la disposición de los electrodos debajo del mentón (músculo Mylohyoid).}
\end{figure}

\begin{figure}[H]
    \centering
    \includegraphics[width=0.45\linewidth]{insumos_ia/cuaderno_de_labo/images/image110.jpg}
    \includegraphics[width=0.45\linewidth]{insumos_ia/cuaderno_de_labo/images/image423.jpg}
    \caption{Vista general del cableado final hacia los tres canales.}
\end{figure}

Llevamos un registro detallado de las pruebas realizadas ese día, evaluando distintas configuraciones y pronunciaciones de palabras. A continuación, se detalla la bitácora de mediciones:

\begin{table}[H]
    \centering
    \caption{Registro de mediciones realizadas el 02/12/2025. FS = Full Scale, GS = Ground Scale.}
    \label{tab:mediciones_labo6}
    \begin{tabularx}{\linewidth}{l l c l X}
        \toprule
        \textbf{Sujeto} & \textbf{Título} & \textbf{Temp.} & \textbf{Estado} & \textbf{Comentario} \\
        \midrule
        Santi & B\_Prueba1\_Santi\_sinfiltronotch & \SI{26}{\degreeCelsius} & Nuevos & Notamos que se despegó un poco el electrodo del mylohyoid. Todos los del NIDAQ en GS. \\
        '' & B\_Prueba2\_Santi\_AI3FS-conNotch & '' & '' & Agregamos el filtro notch y pusimos el AI3 en FS, el resto en GS. \\
        '' & B\_Prueba3\_Santi\_ALLGS-conNotch & '' & '' & Todos los canales en GS, con filtro Notch. \\
        '' & B\_Prueba4\_Santi\_ALLFS-conNotch & '' & semana pasada & Usamos todos en FS con filtro Notch. \\
        '' & B\_Prueba5\_Santi\_ALLFS-conNotch-conFiltro20708 & '' & este dia & Todos en FS con filtro notch y filtro pasabanda: \SI{20}{\hertz} y \SI{708}{\hertz}. \\
        '' & NO\_Prueba1\_Santi\_ALLFS-conNotch & '' & '' & Todos en FS con filtro notch. Palabra \enquote{NO}. Cambiamos el BPM del metrónomo a 40. \\
        '' & JERJES\_Prueba1\_Santi\_ALLFS-conNotch & '' & '' & Cambiamos la palabra a \enquote{JERJES}. \\
        '' & NANDU\_Prueba1\_Santi\_ALLFS-conNotch & '' & '' & La palabra es \enquote{ñandu}, creemos que puede generar problemas la ñ. \\
        '' & ARTE\_Prueba1\_Santi\_ALLFS-conNotch & '' & '' & Notamos que el electrodo inferior del orbicularis oris se habia despegado casi a la mitad. Palabra \enquote{Arte}, última medición. Los músculos parecen estar sincronizados. \\
        '' & NANDU\_Prueba2\_Santi\_ALLFS-conNotch & '' & '' & Segunda prueba de \enquote{ñandu}, con video. \\
        \bottomrule
    \end{tabularx}
\end{table}

Cabe destacar que todas las grabaciones \textbf{se guardaron crudas}, sin el filtro notch aplicado por \textit{hardware} en la placa de adquisición, permitiendo el filtrado y análisis digital de manera \textit{offline}.

Con la experiencia acumulada del análisis de este tipo de señales, se observó que el \enquote{ruido} predominante se correspondía con la frecuencia de red de \SI{50}{\hertz}. Por lo tanto, se propuso aplicar un \textbf{filtro notch} y un \textbf{filtro pasabanda} con frecuencias de corte de \SI{20}{\hertz} y \SI{2000}{\hertz} de manera digital. Esta utilización de filtros permite que el software de adquisición desarrollado se asemeje a los parámetros utilizados nativamente por el software \textit{Spike Recorder} de Backyard Brains.

El desarrollo de una metodología robusta para diferenciar palabras forma parte del Plan de Trabajo. Como cierre del experimento, se analizaron los registros de las palabras \enquote{Arte} y \enquote{Ñandú}, extrayendo su envolvente RMS junto con la señal cruda procesada.

En la palabra \enquote{Arte}, las mediciones presentaron mayor correlación entre sí, observándose además que el músculo \textit{Mylohyoid} (CH0) tuvo una señal de menor magnitud general (Figura \ref{fig:labo6_arte}).

\begin{figure}[H]
    \centering
    \includegraphics[width=\linewidth]{insumos_ia/cuaderno_de_labo/images/plot_calibrado_ARTE_Prueba1_Santi.png}
    \caption{Señales medidas durante la pronunciación de la palabra \enquote{Arte}. Se aprecia la sincronización general y la menor magnitud en el CH0.}
    \label{fig:labo6_arte}
\end{figure}

Por otro lado, en la palabra \enquote{Ñandú} ambas sílabas generaron picos distinguibles en las mediciones. Cuando se pronuncia \enquote{Ñan}, se activan los tres músculos en sincronía, mientras que la sílaba \enquote{dú} activa predominantemente el \textit{Orbicularis oris}, ya que la expresión requiere una contracción labial específica. Además, la señal presenta un perfil más suave (Figura \ref{fig:labo6_nandu}).

\begin{figure}[H]
    \centering
    \includegraphics[width=\linewidth]{insumos_ia/cuaderno_de_labo/images/plot_calibrado_NANDU_Prueba2_Santi.png}
    \caption{Señales medidas durante la pronunciación de la palabra \enquote{Ñandú}. Se distinguen dos picos correspondientes a cada sílaba, con el \textit{Orbicularis oris} dominando el segundo pico.}
    \label{fig:labo6_nandu}
\end{figure}


% --- BIBLIOGRAFÍA ---
\printbibliography

\end{document}
